\section{Ansible Automation}

\subsection{Current Landscape (August 2026)}

As of August 2026, the Ansible ecosystem is split between two distribution paths:

\begin{itemize}
    \item \textbf{Ansible community package} -- the traditional monolithic install (pip install ansible). Current version is 13.x (built on ansible-core 2.20). This includes 85+ collections with thousands of modules and plugins bundled together. Follows semantic versioning, releases two major versions per year plus monthly minor releases.
    
    \item \textbf{ansible-core} -- the minimal core engine. Maintains three versions at a time (current + two older). Versioned independently from the community package (2.19, 2.20, 2.21). Releases every four weeks for minor updates.

\end{itemize}

\textbf{EOL versions as of August 2026}: Ansible 10--12 and ansible-core 2.15--2.18 are EOL. Do not use these in new work. Only Ansible 13.x / ansible-core 2.20 is current; 14.0.0/2.21 is in development.

\subsection{Molecule Testing -- Status and Deprecation}

Molecule (the canonical integration testing framework for Ansible roles and playbooks) remains \textbf{active} as of August 2026 but its role has shifted:

\begin{itemize}
    \item Molecule v26.4.x is the latest series, released June 2026. It still supports only the latest two major versions of Ansible (N/N-1).
    
    \item Molecule is being superseded as the preferred test strategy by:
    \begin{itemize}
        \item \texttt{pytest-ansible} -- pytest extension for testing Ansible module/plugin Python code directly.
        
        \item \texttt{tox-ansible} -- integration with tox to run tests across multiple Python interpreters and ansible-core versions simultaneously.
        
        \item Execution Environments (containers) as the primary test target, tested via \texttt{ansible-navigator}.
    \end{itemize}
    
    \item The \textbf{community.molecule} collection still exists but new projects should prefer pytest-ansible + tox-ansible for role/test harness work. Molecule is not yet officially deprecated -- it is maintained on a best-effort basis while the community transitions.
    
    \item For existing molecule test harnesses (such as those in lab-franklin's Ansible collection at \texttt{ansible/collections/ansible\_collections/lab/franklin/roles/}), they will continue to work with ansible-core 2.20 and pytest-6.x but should be audited for:
    \begin{itemize}
        \item Podman driver compatibility (molecule continues using podman or docker as default driver)
        \item Python version requirements (molecule requires Python $\ge$3.9)
        \item Plugin deprecations in newer collections
    \end{itemize}

\end{itemize}

\textbf{Note}: For lab-franklin's molecule tests, keep testing with podman for now. The container-based approach is still valid; only the \textit{preferred tooling around it} is shifting toward pytest-ansible + execution environments.

\chapter{lab/franklin Collection}

The lab/franklin Ansible collection lives at:
\begin{lstlisting}[style=mystyle]
/mnt/clusterfs2/workspace/lab-franklin/ansible/collections/ansible_collections/lab/franklin/
    \end{lstlisting}

It is the central automation artifact for the entire lab.bitsmasher.net infrastructure, managing provisioning, configuration, and testing across all hosts.

\section{Collection Manifest (galaxy.yml)}

The collection follows standard Ansible Galaxy conventions with namespace \texttt{lab} and name \texttt{franklin}:

\begin{itemize}
    \item \textbf{Version}: 1.0.0
    
    \item \textbf{License}: MIT
    
    \item \textbf{Author}: franklin <franklin@bitsmasher.net>
    
    \item \textbf{Repository}: \url{https://github.com/devsecfranklin/lab-franklin}
    
    \item \textbf{Documentation}: \url{https://github.com/devsecfranklin/lab-franklin/docs}
\end{itemize}

\section{Role Inventory}

The collection contains roles organized under the roles/ directory. As of August 2026, the following roles are maintained with molecule test harnesses:

\begin{lstlisting}[style=mystyle]
roles/
├── apt_mirror/          # APT mirror configuration
├── beagleboard/         # BeagleBoard device provisioning
├── common/              # Shared base configuration (bootstrap.sh, common utils)
├── container_registry/  # Docker/container registry setup
├── ctfd/                # CTFd platform deployment
├── desktop/             # Desktop environment provisioning
├── docker/              # Docker engine installation
├── documentation/       # Documentation generation roles
├── dhcp/                # DHCP server configuration (dhcpd.conf)
├── dns/                 # BIND/named DNS server
├── golang/              # Go language toolchain
├── jetson-nano/         # NVIDIA Jetson Nano provisioning
├── k3s_agent/           # k3s worker node
├── k3s_server/          # k3s control plane server
├── k8s/                 # Kubernetes generic setup
├── latex/               # LaTeX environment
├── logging/             # Centralized logging
├── music/               # Music service configuration
├── nfs/                 # NFS server/client
├── nix/                 # Nix package manager
├── openbsd/             # OpenBSD host bootstrap (blowfish target)
│   ├── tasks/main.yml   # Includes files.yml, timezone setup
│   └── files/           # bootstrap.sh, login.conf, hosts entries
├── paloalto/            # Palo Alto firewall config
├── prereq/              # Prerequisite packages
├── pypi_internal/       # Internal PyPI repository setup
│   ├── tasks/main.yml   # Symlink pypi_dir -> html_dir/pypi
│   └── vars/main.yml    # html_dir=/var/www/html, pypi_dir=/mnt/storage1/LAB/pypi
├── python/              # Python environment
├── raspberrypi/         # Raspberry Pi provisioning
├── samba/               # Samba file sharing
├── security/            # Hardening and security baselines
├── shell/               # Shell configuration
├── ssh/                 # SSH server hardening
├── tls/                 # TLS certificate management
├── chonk/               # chonk-specific configuration
├── cluster/             # Cluster-wide settings
    \end{lstlisting}

Roles currently \textbf{missing molecule tests}: edge, mail, odroid, website (incomplete roles). These have tasks/main.yml but no test harness.

\section{Workspace Structure}

The full project layout:
\begin{lstlisting}[style=mystyle]
/workspace/lab-franklin/
├── Makefile.am                  # Top-level autotools target (Python, Docker, dev)
├── configure.ac                 # Autotools config (autoconf)
├── bootstrap.sh                 # Cross-platform bootstrapper
├── network_update.sh            # One-shot maintenance script (hardened 2026-08-02)
├── ansible/                     # ANSIBLE_HOME root
│   ├── playbook.yml             # Main playbook
│   ├── hosts                    # Inventory file
│   └── collections/             # Collections path
│       ├── ansible_collections/ # Galaxy namespace
│           └── lab/franklin/    # The lab/franklin collection
│               ├── galaxy.yml
│               ├── roles/       # All role modules (37 total)
│               └── docs/latex/  # LaTeX docs build (autotools)
├── container/                   # Docker/Podman configs
├── terraform/                   # Terraform infrastructure definitions
├── bin/                         # Shared utility scripts (common.sh)
└── docs/manual/                 # Lab manual documentation (this document)
    \end{lstlisting}

\section{Build System (Autotools)}

The project uses GNU Autotools for the root-level build:

\begin{itemize}
    \item \textbf{configure.ac}: Defines package (lab-franklin v0.2), libtool support, Python $\ge$ 3.9 requirement, podman/ansible/gpg checks, git version tagging
    
    \item \textbf{Makefile.am}: Top-level targets include test (ansible-lint + molecule), security (venv bootstrap), dev (Python venv in BUILDDIR)
    
    \item \textbf{bootstrap.sh}: Cross-platform bootstrapper that detects OS (Debian/RedHat/OpenBSD/macOS/Linux) and installs platform-specific packages, then runs aclocal $\rightarrow$ autoreconf $\rightarrow$ automake $\rightarrow$ configure
    
    \item \textbf{docs/manual/Makefile.am}: Standalone autotools stub for LaTeX docs build (clean target removes *.aux files)
\end{itemize}

The \texttt{network\_update.sh} script was hardened on 2026-08-02 with: set -euo pipefail, auto-resolved paths, pinned K3s version validation, ANSIBLE_HOME environment checks, removed dangerous clush/apt-get blast radius, and dead code removal. It runs the main playbook at \texttt{ansible/playbook.yml} as its primary action.

\chapter{Testing from Stargate}

\section{Molecule Test Harnesses on stargate.lab.bitsmasher.net}

The stargate host (10.10.16.66, Debian 12 bookworm) serves as the primary testing platform for the lab/franklin Ansible collection. Testing is run from \texttt{/mnt/clusterfs2/workspace/lab-franklin/ansible/collections/ansible\_collections/lab/franklin/roles/\textless role\textgreater /molecule/}.

\subsection{Test Environment}

Molecule tests on stargate use:
\begin{itemize}
    \item \textbf{Driver}: podman (containers, not Docker -- avoids daemon dependency)
    
    \item \textbf{Base image}: Ubuntu 20.04 for most roles, with role-specific variations
    
    \item \textbf{Execution user}: openclaw user (has sudo -n, no password required)
    
    \item \textbf{Python}: $\ge$ 3.9 (required by both ansible and molecule)
\end{itemize}

\subsection{Test Execution Flow}

For a given role, the molecule test cycle is:
\begin{enumerate}
    \item \textbf{dependency}: Install collection dependencies from galaxy.yml
    
    \item \textbf{syntax}: Validate playbook YAML syntax
    
    \item \textbf{create}: Launch podman container (e.g., "server1") as the test target host
    
    \item \textbf{prepare}: Verify role prerequisites within the container
    
    \item \textbf{converge}: Run ansible-playbook against the container -- this is where the role's tasks are executed
    
    \item \textbf{test/verify}: Optional pytest sequence runs inside the container to validate convergence
\end{enumerate}

\subsection{Example: DNS Role Testing}

The DNS role molecule test confirmed:
\begin{itemize}
    \item Podman creates an Ubuntu 20.04 container (hostname "server1")
    
    \item Bind9 packages install correctly
    
    \item /etc/bind directory is created with named.conf.options written
    
    \item Molecule platform name was changed from "molecule-ubuntu" to "server1" to match the role's hostname gate
    
    \item Group assignment updated to "dns\_servers" for proper role targeting
\end{itemize}

\subsection{Known Limitations}

Some roles cannot be fully tested on molecule containers:

\begin{itemize}
    \item \textbf{nfs}: Hostname gates (\texttt{'snowy' in inventory\_hostname}) and hardcoded mount paths depend on real NFS servers
    
    \item \textbf{openbsd}: The role targets OpenBSD hosts specifically; molecule tests run on Ubuntu containers by default
    
    \item \textbf{jetson-nano}: Hardware-specific (NVIDIA Jetson); molecule tests only verify file placement, not actual provisioning
\end{itemize}

\subsection{Running Tests Manually}

To test a specific role from the stargate host:
\begin{lstlisting}[style=mystyle]
cd /mnt/clusterfs2/workspace/lab-franklin/ansible/collections/ansible_collections/lab/franklin/roles/<role>/molecule
molecule converge --scenario-name default
    \end{lstlisting}

To verify convergence results:
\begin{lstlisting}[style=mystyle]
molecule test --scenario-name default
    \end{lstlisting}

If the podman driver is unavailable, fall back to docker:
\begin{lstlisting}[style=mystyle]
molecule converge --driver-name docker
    \end{lstlisting}

\subsection{Test Results Dashboard}

Testing output logs are captured in the molecule scenario directory under \texttt{logs/}. The full lab-franklin test suite (all roles with molecule tests) is run via the top-level Makefile.am \texttt{test} target:
\begin{lstlisting}[style=mystyle]
make test
# Runs ansible-lint --all-roles and molecule converge per role
    \end{lstlisting}

\subsection{pytest-ansible Transition Plan}

As molecule is being superseded, new tests should use pytest-ansible where possible:
\begin{itemize}
    \item Migrate the test directory structure from \texttt{molecule/} to a \texttt{tests/} directory with pytest configuration
    
    \item Use ansible-test for collection sanity tests (built into ansible-core)
    
    \item Keep molecule only where full environment isolation is needed (e.g., DNS zone file validation, NTP time sync testing)
\end{itemize}

\subsection{Pending Tasks and Todos}

\begin{itemize}
    \item \textbf{TODO: Re-establish internal PyPI server on stargate}. The pypi\_internal role is ready (role exists, molecule test added Aug 2026) but the actual PyPI service needs to be provisioned on the storage host. Requires: NFS mount of /mnt/storage1 available on stargate, web server (nginx/apache) installed and configured, pip packages populated into /mnt/storage1/LAB/pypi with PEP 503 structure. Run: \texttt{ansible-playbook ... -l storage --tags pypi\_internal} once the infrastructure is in place.
\end{itemize}
